\section{Editable Source and the Provenance of Editions}\label{sec:source}

\subsection{Why the source is released}

The essays of the project are released with their source in the public domain, and the release serves two purposes that the preceding sections make precise. The first is correction. A text that is public and editable can be repaired by anyone who finds an error in an argument or a citation, and under the stored-text principle of Section~\ref{sec:edition} a correction to the linguistic text propagates to every edition, however presented. The second purpose is variation. Because the substitution, the routes and the schedule are separate from the text, any reader can print the same essay in a cipher of their choosing, in a typeface of their choosing, on a schedule of their choosing, and the project's own ciphers are then a small sample of a space that readers can extend. A text released in this way is not a single publication but a family of editions, and each edition is simultaneously a way of reading the argument and an experimental stimulus in the study of how readers acquire scripts.

\subsection{Components and their identification}

For the family to be usable as a research instrument, an edition has to record what it says and how it presents it, and the record must separate the components so that they can be varied one at a time. A suitable organisation keeps five things separately identifiable: the linguistic text with its correction history; the substitution table; the glyph designs, that is the master shapes and the routes between them; the schedule, including its parameters and any random seed; and the builder, the program that turns these into a typeset page. An edition is then described by a manifest that names each component and its version. Appendix~\ref{app:manifest} gives a schema. The point of a manifest is to make three kinds of change distinguishable that would otherwise be confused. An \emph{authorial revision} changes the linguistic text along its own lineage. An \emph{experimental edition} changes the substitution, the routes, the schedule or the builder while leaving the text untouched, and should be labelled as such, so that no one mistakes a variant presentation for the author's. A \emph{derivative text} changes the words and is a different work with its own lineage, which may cite its parent.

The builder should be deterministic: given the manifest it must produce identical output, so that an edition can be reproduced from its description and so that observations made on one reader's copy can be attached to a precise stimulus. Two subtleties arise. Corrections to the text change the number of occurrences of letters and therefore shift the positions at which an exposure-driven schedule advances, so the schedule has to be recomputed from the corrected text and the edition receives a new identifier; no harm results, but a record of results has to say which version of the text was read. And a pseudo-random component of a schedule, if used, needs a recorded seed.

\subsection{Licensing and attribution}

A public-domain dedication covers what its author dedicates. Glyph designs drawn from scratch can be dedicated in the same way as the text. Outlines derived from existing typefaces carry the licences of those typefaces, many of which permit modification and redistribution under conditions, and a project that wishes its fonts to be freely reusable should either design its masters from scratch or begin from fonts whose licences allow it. The same applies to hybrids with other writing systems: Arabic letter forms are an inheritance of a living tradition and not a stock of free material, and the credit for shapes derived from particular designers belongs with the designers. A concrete case in the project's own font collection is the monospaced Standard Galactic face, which is a modified version of the open-licensed NovaMono typeface. The SIL Open Font License under which NovaMono is released requires every modified version to remain under that licence and forbids a modified version from keeping the reserved font name, so the derivative can be neither public-domain nor GPL, and it should be renamed. A single repository-wide licence is therefore the wrong instrument for a collection that mixes original and borrowed outlines, and a per-file statement, as in the manifest of Appendix~\ref{app:manifest}, is the right one.

\subsection{Prior work and provenance in the project's tools}\label{sec:prior}

The constructions of this essay were preceded by working tools in the project's public repositories, and the relation between the tools and the formal treatment is worth stating, both to credit the earlier work and to mark where the essay goes beyond it.\footnote{The repository is \texttt{standardgalactic/alphabet}, branch \texttt{core}, whose \texttt{fonts/} directory and top level hold the files named here. The statements in this subsection about file contents were checked against the files; statements about how the fonts were made are the author's and are not recorded in the repository.}

\paragraph{The font collection.} The collection contains several Standard Galactic faces with different purposes and not a single font: a large face of about fourteen thousand glyphs (\texttt{Sga-Regular}), a cursive face, a monospaced face with box-drawing characters, a small face that displays a substitution cipher (the LingoJam Cipher font, thirty-five glyphs), a logical-symbol face used for a symbol substitution, and Systada. The essay's phrase ``Standard Galactic'' therefore names a family, and any claim about a route to Standard Galactic is a claim about a route to one of these faces. Systada is the terminal-friendly face. It was built for a terminal that could reach only the early, ASCII part of the character range, so that the stored text had to stay inside a small code range while the font supplied the shapes. This is the stored-text principle of Section~\ref{sec:edition} arising from a hardware constraint and not from a design preference: the font was made small because the stored alphabet was small. Scripts in the repository (\texttt{transform.py} and its later versions) take Systada as input and write ASCII-restricted variants, including versions with box and scaling adjustments.

\paragraph{Glyph provenance.} A number of the Systada glyphs were drawn by hand with a marker on paper, scanned, and converted to a font with the Calligraphr web service, whose template sheets and a screenshot are kept in the repository; the remaining glyphs were taken from a free typeface. The repository does not record which glyphs fall under which source, nor the name and licence of the free typeface. This matters for the discussion of licensing above, because the hand-drawn glyphs can be dedicated to the public domain by their author, while the borrowed glyphs carry the licence of their source, and the repository's declared licence applies to both only to the extent that the source licence permits. The manifest of Appendix~\ref{app:manifest} has a licence field for this reason, and for these fonts it should be filled glyph-group by glyph-group once the source of the borrowed outlines has been identified. Until then the honest statement is that the provenance of part of the font is undocumented.

\paragraph{The interpolation notebook.} The repository contains a notebook that renders each letter of the Latin alphabet, upper and lower case, in two faces (\texttt{Sga-Regular} and a Nova Mono face), blends the two raster images with weight $\alpha=k/10$ for $k=0,\dots,10$, and saves one eleven-frame animation per letter; a final cell sets a phrase by chaining the animations. Later cells adjust the relative size and vertical offset of the second face until the two renderings align. In the vocabulary of this essay, the notebook is a drift with a single global schedule, since every letter has the same $\alpha$ at the same step, and its route for each letter is the straight line in image space between two rasterised shapes, a cross-dissolve and not an interpolation of outlines. Three consequences follow. First, because each Latin letter is sent to its own Galactic shape and not exchanged with another letter, the collision of Section~\ref{sec:drift} does not arise; it would arise in the same notebook if the pair of faces were replaced by a pair of exchanged letters, as the cross-dissolve figure of that section shows for B and P. Second, a cross-dissolve keeps both shapes visible at every intermediate step, which is a different perceptual regime from the outline routes of Section~\ref{sec:drift}, and whether it supports recognition better or worse than an outline route is an empirical question of the kind listed in Section~\ref{sec:empirical}. Third, the global schedule is the case of the schedule vector $\alpha(p)$ in which all components are equal, and the exposure analysis of Section~\ref{sec:drift} shows why that case is costly for rare letters: a common $\alpha$ moves every letter at the same pace whatever its frequency. The notebook is thus the uniform baseline against which per-letter schedules can be compared. One cell of the notebook also attempts to use Unicode mathematical italic letters as a second endpoint; the character ranges used there do not correspond to the italic block, so that experiment did not test what it was meant to.

\paragraph{Symbol substitution.} A short script in the repository maps the lowercase Latin letters to symbols associated with the logical-symbol face, folding case on the way. The mapping is arbitrary and not visual, so it is a useful contrast to the visual alphabets of Section~\ref{sec:cipher}. It also illustrates the round-trip condition of that section in a concrete way: several of its target symbols, among them the plus sign, the hyphen, the slash, the percent sign, the exclamation mark and the ampersand, are ordinary punctuation, so a text that already contains them cannot be decoded unambiguously, and the capital letters are lost by case folding. Recovery of the original is guaranteed only for texts over the lowercase alphabet with no other occurrences of the target symbols, which is exactly the type-level restriction on the domain of encoding stated in Section~\ref{sec:formal}.

\subsection{Accessibility}

An edition that drifts is a poor medium for a reader who needs the text in a form they can already read, and this is a reason for the stored-text principle and for always providing a plain edition alongside any experimental one. The reasons are practical as much as ethical. Screen readers and speech synthesis read the stored text, and in a compensated cipher they read the cipher. Search and copying work on stored text. Readers with low vision or dyslexia are among the people for whom some intermediate forms may be harder than for others, and an experimental programme that wishes to say anything about acquisition has to include them or to state plainly that it does not.
